\documentclass{situationist-preprint}

\title{Wittgenstein and Davidson Reconciled}
\msid{LNG-01-b}
\msfamily{LNG-01}
\msdiscipline{Philosophy of Language}
\mstype{Companion manuscript}
\mskeywords{meaning as use; truth-conditional semantics; Wittgenstein; Davidson; radical interpretation}

\begin{document}
\maketitle

\begin{abstract}
Wittgenstein's remark that the meaning of a word is, in most cases, its use in the language and Davidson's proposal that a theory of meaning takes the form of a theory of truth are usually presented as the two poles of the subject. LNG-01-a argues that meaning is use constrained by truth-conditions. We argue that this is not a compromise between the two authors but a statement each of them, read carefully, is committed to. Davidson's truth theory is answerable to use through radical interpretation; Wittgenstein's use is answerable to truth through the practice of assessing applications as correct. The residual disagreement concerns the order of explanation, and the Reconciliation theorem of LNG-01-a settles it.
\end{abstract}

\section{Two theses and their commitments}

Wittgenstein's thesis is that what a word means is fixed by the way it is used in the practice of a linguistic community \citep[§43]{Wittgenstein1953}. Davidson's is that to give the meaning of a sentence is to state its truth-condition in a finite, compositional theory of truth for the language \citep{Davidson1967}. Read as rival definitions, they disagree about what meaning is. Read as theses with commitments, they converge.

Davidson's theory is empirical. A truth theory is correct for a speaker's language only if it is the theory a radical interpreter would arrive at from the speaker's dispositions to hold sentences true in observable circumstances \citep{Davidson1973}. What the interpreter has to go on is the speaker's use: the circumstances in which she applies her sentences and the transitions she makes between them. A truth theory that does not fit that use is not a theory of the speaker's meaning, however elegant. So Davidson accepts that use constrains truth-conditions; in the terms of LNG-01-a, an assignment counts only if it is sound for the rules the speaker follows (Definitions~2.3 and~2.4 there).

Wittgenstein's use is not a mere pattern of behavior. It is a practice of following rules, in which applications are assessed as correct or incorrect \citep[§§199--202]{Wittgenstein1953}. For declarative sentences, the assessment that matters is whether the sentence was correctly asserted, and in the central cases a sentence is correctly asserted when what it says is so. A use whose rules lead from correct assertions to incorrect ones is defective by the practice's own lights. Wittgenstein is therefore committed to the truth-conditional constraint, Principle~P2 of LNG-01-a: a use counts as conferring meaning only if its rules are truth-preserving under some assignment.

\section{The residual disagreement}

What remains is the order of explanation. For Davidson, truth is primitive and meaning is explained through it; for Wittgenstein, use is primitive and truth-talk is one move within it. The Reconciliation theorem (Theorem~3.6 of LNG-01-a) shows that the choice does not affect what a given expression means. Given Principles~P2 and~P3, an expression contributes the value that every assignment admissible for its use gives it, and nothing if no assignment is admissible. Use is what is constrained; truth-conditions are what constrains it; and the value contributed is the same whichever is called basic. Where the use is determinate, as Theorem~3.3 shows for the standard connectives, a Davidsonian truth theory and a Wittgensteinian description of the practice issue in the same clause for each expression.

The theorem also explains the one case in which the two authors seem to diverge. Prior's tonk has a use, in the sense that its rules can be stated and followed, and a Wittgensteinian might be expected to grant it a meaning; a Davidsonian has no truth clause to give it. Proposition~3.2 and Corollary~3.7 of LNG-01-a show that tonk contributes nothing. A Wittgensteinian who holds, as above, that a practice's rules must preserve correctness of assertion agrees: the tonk practice is not a language game in which anything is said.

\begin{objection}
Wittgenstein rejected systematic theories of meaning, and Davidson's programme is nothing if not systematic. The reconciliation ignores this.
\end{objection}

\begin{reply}
The reconciliation concerns what meaning is, not whether the meanings of a natural language can be exhaustively codified. The account of LNG-01-a requires only that some compositional assignment be admissible for the use, not that speakers or theorists be able to state it. Wittgenstein's suspicion of theory is directed at the idea that a rule's application is fixed by something behind the practice; Davidson's truth theory, constrained by radical interpretation, places nothing behind it.
\end{reply}

\section{Conclusion}

Davidson's truth theory answers to use, and Wittgenstein's use answers to truth. Each author is committed to the thesis that meaning is use constrained by truth-conditions, and the Reconciliation theorem shows that their disagreement about which notion is basic leaves the meaning of every expression unchanged.

\begin{thebibliography}{9}\small
\bibitem[Davidson(1967)]{Davidson1967} Davidson, D. (1967). Truth and meaning. \emph{Synthese}, 17, 304--323.
\bibitem[Davidson(1973)]{Davidson1973} Davidson, D. (1973). Radical interpretation. \emph{Dialectica}, 27, 313--328.
\bibitem[Wittgenstein(1953)]{Wittgenstein1953} Wittgenstein, L. (1953). \emph{Philosophical Investigations}. Translated by G.~E.~M. Anscombe. Blackwell.
\end{thebibliography}

\end{document}
